%%=============================================================================
%% Samenvatting
%%=============================================================================

% TODO: De "abstract" of samenvatting is een kernachtige (~ 1 blz. voor een
% thesis) synthese van het document.
%
% Een goede abstract biedt een kernachtig antwoord op volgende vragen:
%
% 1. Waarover gaat de bachelorproef?
% 2. Waarom heb je er over geschreven?
% 3. Hoe heb je het onderzoek uitgevoerd?
% 4. Wat waren de resultaten? Wat blijkt uit je onderzoek?
% 5. Wat betekenen je resultaten? Wat is de relevantie voor het werkveld?
%
% Daarom bestaat een abstract uit volgende componenten:
%
% - inleiding + kaderen thema
% - probleemstelling
% - (centrale) onderzoeksvraag
% - onderzoeksdoelstelling
% - methodologie
% - resultaten (beperk tot de belangrijkste, relevant voor de onderzoeksvraag)
% - conclusies, aanbevelingen, beperkingen
%
% LET OP! Een samenvatting is GEEN voorwoord!

%%---------- Nederlandse samenvatting -----------------------------------------
%
% TODO: Als je je bachelorproef in het Engels schrijft, moet je eerst een
% Nederlandse samenvatting invoegen. Haal daarvoor onderstaande code uit
% commentaar.
% Wie zijn bachelorproef in het Nederlands schrijft, kan dit negeren, de inhoud
% wordt niet in het document ingevoegd.

\IfLanguageName{english}{%
\selectlanguage{dutch}
\chapter*{Samenvatting}
\lipsum[1-4]
\selectlanguage{english}
}{}

%%---------- Samenvatting -----------------------------------------------------
% De samenvatting in de hoofdtaal van het document

\chapter*{\IfLanguageName{dutch}{Samenvatting}{Abstract}}

Het Virtual IT Company (VIC) van HOGENT maakt binnen zijn Proxmox-omgeving momenteel uit sluitend gebruik van klassieke virtual machines (VM’s). Dit onderzoek analyseert de potentiële meerwaarde van LXC-containers als efficiënter en lichter alternatief voor specifieke workloads. De focus ligt hierbij op performance, beveiliging, beheerbaarheid en de verschillen tussen Proxmox versie 8 en 9. Via een literatuurstudie en een comparatieve proof of concept worden de technische eigenschappen van beide virtualisatietechnologieën tegen elkaar afgezet. Hierbij wordt onderzocht in welke scenario’s de lagere over-head van containers opweegt tegen de striktere isolatie van VM’s. Het eindresultaat is een adviesrapport met beslissingscriteria, waarmee IT-beheerders fundamentele keuzes kunnen maken voor de inrichting van de infrastructuur en toekomstige ontwikkelingen zoals Kubernetes.

%% TODO: Werk de samenvatting bij na afronding van de PoC met de definitieve resultaten en conclusies.
